% ---------------------------------------------------------------------------
%  paper/shared/disclosure.tex — the two authorship disclosure blocks.
%
%  docs/publication-split.md §1 fixes the rule these implement:
%
%    "Under no circumstances will there be any ambiguity about whether a paper
%     is wholly my work (none), merely my prose (some) or entirely LLM-authored
%     (some)."   — jasonp, 2026-08-07
%
%  There are exactly two blocks and they live in exactly one file, so that a
%  paper cannot quietly soften its own disclosure.  The fixed sentences take no
%  arguments.  The ONE thing a paper supplies is the per-result verification
%  ledger, because that is the part that differs honestly between papers — and
%  §1 requires it to be per result, not in aggregate.
%
%  Do not add a \Pdisclosure or \Ldisclosure variant with an optional softening
%  argument.  If a paper does not fit one of these two blocks, the paper is
%  wrong, not the block.
% ---------------------------------------------------------------------------

\newsavebox{\disclosurebox}

\newenvironment{disclosureframe}
  {\begin{center}\begin{lrbox}{\disclosurebox}\begin{minipage}{0.93\textwidth}\small}
  {\end{minipage}\end{lrbox}\fbox{\usebox{\disclosurebox}}\end{center}}

% --- P: jasonp's prose -------------------------------------------------------
% Byline: Jason Parker, alone.  Argument: the per-result ledger of what the
% machine did.
\newcommand{\Pdisclosure}[1]{%
\begin{disclosureframe}
\textbf{Authorship.}\enspace Every sentence of this paper was written by its
named author. He has read every proof it states and believes each one, and he
is responsible for the correctness of what it claims.

\smallskip
\textbf{Machine contribution.}\enspace #1
\end{disclosureframe}}

% --- L: entirely machine-written --------------------------------------------
% Argument: the per-result verification ledger — for each result, what warrants
% it (Lean kernel, exact certificate, or labelled measurement) and how much of
% it a human has checked.
\newcommand{\Ldisclosure}[1]{%
\begin{disclosureframe}
\textbf{Authorship.}\enspace This document was written end to end by a large
language model. That includes its mathematics, its prose, its tables and its
\LaTeX{}. No sentence of it was written by a human. The mathematics it reports
was likewise derived by machine.

\smallskip
\textbf{Human role.}\enspace The person who directed this work chose what to
compute, chose what to publish, and is responsible for the decision to release
this document. Except where the ledger below says otherwise, that person has
not personally checked the arguments here, and this disclosure is not to be
read as an endorsement of them.

\smallskip
\textbf{What warrants each result.}\enspace #1
\end{disclosureframe}}

% --- Draft banner ------------------------------------------------------------
% For an L-paper whose verification ledger is not yet settled.  Loud on purpose:
% an L-paper without a truthful ledger must not be mistaken for a finished one.
%
% \firstdraftbanner is the standard wording, which five papers previously
% carried character-for-character in their own preambles.  It lives here for the
% same reason the disclosure blocks do: identical text in seven places drifts,
% and a paper must not be able to soften it by editing its own copy.  The
% optional argument is for a gate specific to one paper, appended verbatim; a
% paper needing different wording (L6, compute-gated) calls \draftbanner direct.
\newcommand{\firstdraftbanner}[1][]{%
\draftbanner{This is a first draft. The verification ledger below records what
warrants each result \emph{mechanically}; the column that records what a human
has read is empty, deliberately, because nobody has read it yet. Do not
circulate until that column says something true.#1}}

\newcommand{\draftbanner}[1]{%
\begin{center}
\fcolorbox{red!60!black}{yellow!15}{%
  \begin{minipage}{0.93\textwidth}\small
  \textbf{DRAFT — NOT FOR CIRCULATION.}\enspace #1
  \end{minipage}}
\end{center}}
